% Reference copy of Section 4.1 before the reader-first rewrite.
% This file is not included in the manuscript.

\subsection{Signed monomial dual}
\label{subsec:signed-set-process}

For~$A\Subset\Lambda$, $P_t\chi_A(\eta)$ is the expected product of the spins
in~$A$ at time~$t$ when the process starts from~$\eta$. The dual runs backward
from~$A$ and records the sites at earlier times through which these spins depend
on the past. We first define a general signed version of this set-valued process.
The death, split, and birth notation below is convenient for the patch construction {\color{red} Maybe a reference to standard literature, that this is the same duality, with a note we just use more convenient notation for us}.
We then choose its rates and signs from the generator~$\Lcal$.

A signed active set is a pair
\begin{equation*}
Y=(A,\sigma)
\in
\cb{A:A\Subset\Lambda}\times\cb{+,-}.
\end{equation*}
We call~$A$ the active set. Along the dual, it records the sites reached by
tracing the dependence backward, while~$\sigma$ records the accumulated sign.
Signs are multiplied in the usual way and act on real
numbers by~$+x=x$ and~$-x=-x$. Fix, for now, rates and sign labels
\begin{equation*}
\delta_i(S),\ \beta_i(S)\geq0,
\qquad
\sigma_i^\delta(S),\ \sigma_i^\beta(S)\in\cb{+,-},
\qquad
i\in\Lambda,\quad S\subseteq N(i).
\end{equation*}
We use the convention~$\beta_i(\varnothing)=0$. For~$i\in A$, define
\begin{align*}
D_{i,S}(A,\sigma)
&=
\bb{\bb{A\setminus\cb{i}}\cup S,
\sigma\sigma_i^\delta(S)},
\\
B_{i,S}(A,\sigma)
&=
\bb{A\cup S,
\sigma\sigma_i^\beta(S)}.
\end{align*}
These are the usual additive set-valued transitions in a form convenient here {\color{red} Usual in what sense? The reader doesnt know.}.
The update~$D_{i,\varnothing}$ is a \emph{death}, which removes~$i$. For
$S\neq\varnothing$, the update~$D_{i,S}$ is a \emph{split}, which replaces~$i$
by~$S$, and~$B_{i,S}$ is a \emph{birth}, which keeps~$i$ and adds~$S$.
Ignoring the sign coordinate, these updates preserve unions, which is the
standard notion of additivity for set-valued particle systems
\cite{Harris1978,Griffeath1979,Gray1986}. The resulting signed set-valued
process has generator
\begin{align*}
\Dcal g(A,\sigma)
={}&
\sum_{i\in A}\sum_{S\subseteq N(i)}
\delta_i(S)
\bb{g\bb{D_{i,S}(A,\sigma)}-g(A,\sigma)}
\\
&+
\sum_{i\in A}
\sum_{\substack{S\subseteq N(i)\\S\neq\varnothing}}
\beta_i(S)
\bb{g\bb{B_{i,S}(A,\sigma)}-g(A,\sigma)}.
\end{align*}

\label{subsec:rates-to-dual}
For the spin system with generator~$\Lcal$, we now choose these parameters by
matching the monomials in~$\Lcal\chi_A$. For~$r\in\R$, let
$\operatorname{sgn}_{\pm}(r)$ be~$+$ when~$r\geq0$ and~$-$ when~$r<0$. For the
coefficients in~\eqref{eq:rate-expansion}, set
\begin{equation}
\label{eq:signed-dual-coefficients}
a_i^\delta(S)=c_i^0(S),
\qquad
a_i^\beta(S)=-c_i^0(S)-c_i^1(S).
\end{equation}
In the expansion of~$\Lcal\chi_A$, the~$a_i^\delta(S)$ is the coefficient of
the monomial indexed by~$\bb{A\setminus\cb{i}}\cup S$, while~$a_i^\beta(S)$
is the coefficient of the monomial indexed by~$A\cup S$. See
Appendix~\ref{app:generator-action}. We choose the corresponding rates and
signs by
\begin{equation}
\label{eq:death-split-data}
\delta_i(S)=\abs{a_i^\delta(S)},
\qquad
\sigma_i^\delta(S)=\operatorname{sgn}_{\pm}\bb{a_i^\delta(S)}
\end{equation}
for every~$S\subseteq N(i)$, and
\begin{equation}
\label{eq:birth-data}
\beta_i(S)=\abs{a_i^\beta(S)},
\qquad
\sigma_i^\beta(S)=\operatorname{sgn}_{\pm}\bb{a_i^\beta(S)}
\end{equation}
for~$S\neq\varnothing$. As above,~$\beta_i(\varnothing)=0$, and we set
$\sigma_i^\beta(\varnothing)=+$.

The coefficient~$a_i^\beta(\varnothing)$ corresponds to an empty-target birth,
which leaves the active set unchanged. It is included in the potential rather
than represented as a jump. Define
\begin{equation}
\label{eq:dual-potential}
V(A)=\sum_{i\in A}V_i,
\qquad
V_i=
\sum_{S\subseteq N(i)}\delta_i(S)
+
\sum_{\substack{S\subseteq N(i)\\S\neq\varnothing}}\beta_i(S)
+
a_i^\beta(\varnothing).
\end{equation}
Since the neighbourhood sizes and flip rates are uniformly bounded, so are the
total dual jump rate per active site and~$\abs{V_i}$. Each jump activates at
most~$\sup_i\abs{N(i)}$ new sites.

\paragraph{Feynman--Kac duality.}
\label{subsec:monomial-duality}
{\color{red} Remove paragraph headers}
For~$Y=(A,\sigma)$, define
\begin{equation}
\label{eq:monomial-duality-function}
H(Y,\eta)=\sigma\chi_A(\eta).
\end{equation}
We write~$\pP_A$ and~$\E_A$ for the law and expectation of the signed set
process started from~$(A,+)$.

\begin{theorem}[Monomial Feynman--Kac duality]
\label{thm:monomial-duality}
For every signed active set~$Y=(A,\sigma)$,
\begin{equation}
\label{eq:monomial-generator-duality}
\Lcal_\eta H(Y,\eta)
=
\Dcal H(Y,\eta)+V(A)H(Y,\eta).
\end{equation}
Consequently, for every~$A\Subset\Lambda$,~$t\geq0$, and~$\eta\in\Omega$,
\begin{equation}
\label{eq:monomial-feynman-kac-duality}
P_t\chi_A(\eta)
=
\E_A\Cb{
\sigma_t
\exp\bb{\int_0^t V(A_s)\,ds}
\chi_{A_t}(\eta)
}.
\end{equation}
\end{theorem}

The coefficient calculation proving~\eqref{eq:monomial-generator-duality} and
the finite-volume passage to~\eqref{eq:monomial-feynman-kac-duality} are given
in Appendix~\ref{app:dual-calculation} {\color{red} Also a reference to relevant literature should be made, duality is standard, this is not new}. Condition~\eqref{eq:fk-integrability}
supplies the domination in the infinite-volume passage {\color{red} Unclear to the reader what this means because it references proof method, which is not important here because duality is standard, and reader probably will skip appendix anyway. Just say the representation holds under integrability condition (FK)}. The pure-death noise
generator acts diagonally on monomials:
\begin{equation*}
\Ncal^{\mb 0}\chi_A=-\abs A\chi_A.
\end{equation*}
Adding~$\varepsilon\Ncal^{\mb 0}$ leaves the dual jumps unchanged and replaces
$V(A)$ by~$V(A)-\varepsilon\abs A$. We therefore call these independent
$1\to0$ transitions \emph{pure deaths}.
